\makeatletter
\graphicspath{%
{./figs/}
{./common/rawfigs/}
{./common/ressources/gradient_factors/construction/}
{./common/ressources/gradient_factors/definition/}
{./common/ressources/gradient_factors/effet_high/}
{./common/ressources/gradient_factors/effet_low/}
{./common/ressources/gradient_factors/recommandation/}
{./common/ressources/detendeur/}
}
\setbeamertemplate{footline}
{%
  \leavevmode%
  \hbox{\begin{beamercolorbox}[wd=.5\paperwidth,ht=2.5ex,dp=1.125ex,leftskip=.3cm plus1fill,rightskip=.3cm]{author in head/foot}%
    \usebeamerfont{author in head/foot}\insertshortauthor
  \end{beamercolorbox}%
  \begin{beamercolorbox}[wd=.5\paperwidth,ht=2.5ex,dp=1.125ex,leftskip=.3cm,rightskip=.3cm plus1fil]{title in head/foot}%
    \usebeamerfont{title in head/foot}\insertshorttitle\nobreak\hfill\insertframenumber{} / \inserttotalframenumber\usebeamercolor[fg]{page number in head/foot}\usebeamerfont{page number in head/foot}\usebeamertemplate{page number in head/foot}
  \end{beamercolorbox}}%
  \vskip0pt%
}
\AtBeginSection[]{
  \begin{frame}
  \vfill
  \centering
  \begin{beamercolorbox}[sep=8pt,center,shadow=true,rounded=true]{title}
    \usebeamerfont{title}\insertsectionhead\par%
  \end{beamercolorbox}
  \vfill
  \end{frame}
}
\setbeamercovered{dynamic}
\beamertemplatetransparentcovered
\newlength{\@minione}
\newlength{\@minitwo}
\newcommand{\juxt}[3][0.48]{%
\setlength{\@minione}{#1\textwidth}%
\setlength{\@minitwo}{0.96\textwidth - \@minione}%
\begin{minipage}[c][5cm][c]{\@minione}
\raggedright#2
\end{minipage}\hfill
\begin{minipage}[c][5cm][c]{\@minitwo}
\raggedright#3
\end{minipage}%
}
\newlength\p@length
\newcommand{\danger}[1]{%
\begin{tikzpicture}[baseline=(X.center)]
\path[use as bounding box] (-0.4,-0.4) rectangle ++(0.8,0.8);
\node[draw=red,very thick,text=red,regular polygon, regular polygon sides=3,rounded corners, inner sep=1pt,font=\bf] (X) at (0,0) {!};%
\end{tikzpicture}
\setlength{\p@length}{\textwidth}%
\addtolength{\p@length}{-1cm}%
\hfill\parbox{\p@length}{\raggedright#1}}

\newcommand{\includefullheightgraphics}[1]{%
\global\beamer@shrinktrue%
\gdef\beamer@shrinkframebox{%
  \setbox\beamer@framebox=\vbox to 0.96\beamer@frametextheight{%
     \centering%
     \includegraphics[height=0.96\beamer@frametextheight]{#1}%
 }%
}}

%%%%%%%%%%%%%%%%%%%%%%%
%%% TOC
%%%%%%%%%%%%%%%%%%%%%%%

\AtBeginSection[]{%
\begin{frame}
\tableofcontents[
currentsection,
hideothersubsections
]
\end{frame}}

\AtBeginSubsection[]{%
\begin{frame}
\tableofcontents[
currentsection,
currentsubsection,
sectionstyle=show/hide,
subsectionstyle=show/shaded,
subsubsectionstyle=show/show/hide
]
\end{frame}
}

\AtBeginSubsubsection[]{%
\begin{frame}
\tableofcontents[
currentsection,
currentsubsection,
sectionstyle=show/hide,
subsectionstyle=show/shaded,
subsubsectionstyle=show/shaded/hide
]
\end{frame}
}

\newcommand{\shadow@section}[1]{%
  \refstepcounter{section}%
  \global\advance\beamer@tocsectionnumber by 1\relax%
   \long\def\secname{#1}%
   \addtocontents{toc}{\protect\beamer@sectionintoc{\the\c@section}{#1}{\the\c@page}{\the\c@part}%
   {\the\beamer@tocsectionnumber}}%
}%
\newcommand{\zsection}[1]{\shadow@section{#1}}

%%%%%%%%%%%%%%%%%%%%%%%
%%%
%%%%%%%%%%%%%%%%%%%%%%%

%%% to get frame height
% magic number from log
\newlength{\frametitleheight}
\setlength{\frametitleheight}{20.4396pt}
\newcommand{\frameheight}{\textheight - \frametitleheight}

%%%%
\tikzfading[name=fade inside,inner color=transparent!80, outer color=transparent!30]
\tikzset{
 unvisible/.style={opacity=0},
 visible on/.style={alt={#1{}{unvisible}}},
 unvisible on/.style={alt={#1{unvisible}{}}},
 covered/.style={opacity=0.3},
 uncover on/.style={alt={#1{}{covered}}},
 cover on/.style={alt={#1{covered}{}}},
 alt/.code args={<#1>#2#3}{%
   \alt<#1>{\pgfkeysalso{#2}}{\pgfkeysalso{#3}} % \pgfkeysalso doesn't change the path
 },
 bubble/.style={ball color=white,path fading=fade inside},
 bobo/.style={starburst,starburst points=10,starburst point height=3pt,draw=red, fill=yellow,minimum width=#1,minimum height=#1,inner sep=2pt}
}
\pgfdeclarelayer{background}
\pgfsetlayers{background,main}
%%%%%
%%% info
\colorlet{infoc}{green!60!black}
% 1. Define the custom color
\setbeamercolor{infoed text}{fg=infoc}

% 2. Define the environment (must be named 'infoenv')
\newenvironment<>{infoenv}{%
  \only{\setbeamercolor{item projected}{bg=infoc}}
  \begin{altenv}%
    {\usebeamercolor[fg]{infoed text}\usebeamerfont{alerted text}}%
    {\usebeamercolor{normal text}\usebeamerfont{normal text}}%
    {\color{.}}{}%
}{\end{altenv}}

% 3. Define the command (optional, for easier usage like \alert)
%\newcommand<>{\info}{\begin{infoenv}#2\end{infoenv}}

\renewcommand{\tableFactor}{0.94}

%% detendeur
\tikzset{%
ouverture/.style={shorten <=-0.1pt,shorten >=-0.1pt,line width=2mm,#1},
coque/.style={gray,line width=3mm},
membrane/.style={blue,line width=1mm},
ressort/.style={ultra thick,orange,decoration={coil,aspect=0,segment length=4mm,amplitude=4mm},decorate},
ressortdet/.style={ultra thick,orange,decoration={coil,aspect=0,segment length=4.25mm,amplitude=4mm},decorate},
bigressort/.style={line width=1.5mm,orange,decoration={coil,aspect=0,segment length=5mm,amplitude=5mm},decorate},
smallressort/.style={thick,orange,decoration={coil,aspect=0,segment length=2mm,amplitude=4mm},decorate},
tinyressort/.style={orange,decoration={coil,aspect=0,segment length=2mm,amplitude=1mm},decorate},
annot/.style={very thin,white,double=lightgray,line cap=round},
gilet/.style={green,thick,rounded corners=1pt, nearly transparent},
piston/.style={darkgray,line width=2mm},
pistonaxis/.style={piston,line width=1mm},
pistonthinaxis/.style={piston,line width=0.7mm}
}
\tikzfading[name=fade down,
top color=transparent!0,
bottom color=transparent!100]
\tikzfading[name=fade right,
left color=transparent!0,
right color=transparent!100]
\colorlet{eau}{cyan!80!blue}
\colorlet{rhp}{cyan!70}
\colorlet{rpi}{cyan!20}
\colorlet{rpa}{cyan!10}
\colorlet{joint}{green!60!black}
\pgfdeclarelayer{background}
\pgfsetlayers{background,main}

\makeatother
